\section{Заключение}

В работе исследованы методы оптимизации с предобуславливанием и затуханием весов. Основные результаты состоят в следующем.

Во-первых, показано, что затухание весов эквивалентно применению метода с предобуславливанием к изменённой задаче: вместо исходной регуляризованной функции $F = f + r$ метод оптимизирует функцию $\widetilde{F}_t = f + \widetilde{r}_t$ с адаптивным покоординатным регуляризатором $\nabla \widetilde{r}_t = D_t \nabla r$. Доказаны существование и гладкость $\widetilde{r}_t$, а также нижняя оценка расстояния между решениями исходной и изменённой задач, показывающая, что эти решения различаются тем сильнее, чем дальше матрица предобуславливания от единичной.

Во-вторых, доказаны четыре теоремы о скорости сходимости. Теорема \ref{theor:1} даёт оценку $\mathcal{O}(1/T)$ на убывание квадрата нормы градиента изменённой функции в невыпуклом детерминированном случае --- с точностью до порога, контролируемого параметром импульса $\beta$. Теорема \ref{theor:2} устанавливает линейную скорость сходимости к решению изменённой задачи в сильно выпуклом случае для регуляризации Тихонова. Теоремы \ref{theor:3} и \ref{theor:4} переносят оба результата на стохастическую постановку с несмещённым градиентом ограниченной дисперсии: в невыпуклом случае получена скорость $\mathcal{O}(1/\sqrt{T})$, совпадающая по порядку со скоростью стохастического градиентного спуска, а в сильно выпуклом --- линейная сходимость до <<шумового шара>>, радиус которого пропорционален шагу и дисперсии градиента.

В-третьих, проведены вычислительные эксперименты на трёх масштабах --- от задач с точно вычислимыми решениями через ResNet-18 и ViT-Tiny на CIFAR-10/100 до предобучения GPT-2 (124M параметров) и дообучения RoBERTa-base на GLUE, --- иллюстрирующие и количественно подтверждающие теорию: у методов с затуханием весов норма градиента исходной функции стабилизируется на ненулевом уровне, тогда как норма градиента изменённой функции убывает в согласии с доказанными оценками, и это поведение воспроизводится по запускам из разных начальных точек. Показано также, что адаптивная регуляризация автоматически ослабляет штраф для параметров нормировочных слоёв и эмбеддингов (воспроизводя практику ручных списков no-decay), а преимущество отделённого затухания весов в качестве реализуется прежде всего через лучшую достижимую точку оптимизации при сопоставимом обобщении на равном уровне подгонки. Отдельно проанализирован новый, промежуточный способ добавления регуляризации (масштабированное затухание весов, AdamWH/OASISWH): установлены его стационарная точка и порог устойчивости, показано сужение области применимости с ростом масштаба задачи.

Ограничения проведённого анализа задают направления дальнейшей работы: учёт первого момента ($\beta_1 > 0$) и коррекции смещения, присутствующих в практическом Adam; анализ предельного поведения матрицы $D_t$; перенос оценок на неограниченные регуляризаторы в невыпуклом случае; масштабирование экспериментов на модели большего размера и стабилизация масштабированного затухания весов (выбор нижней срезки $\alpha \gtrsim \eta\lambda$).
